% Generated draft: requires mathematical and source review.
\documentclass[11pt]{article}
\usepackage[margin=25mm]{geometry}
\usepackage{fontspec}
\setmainfont{Latin Modern Roman}
\usepackage{amsmath,amssymb,mathtools,unicode-math}
\setmathfont{Latin Modern Math}
\usepackage{enumitem,needspace,xurl,hyperref}
\setlength{\parindent}{0pt}
\setlength{\parskip}{5pt}
\setlength{\emergencystretch}{3em}
\newcommand{\N}{\mathbb{N}}
\newcommand{\Z}{\mathbb{Z}}
\newcommand{\Q}{\mathbb{Q}}
\newcommand{\R}{\mathbb{R}}
\newcommand{\C}{\mathbb{C}}
\newcommand{\F}{\mathbb{F}}
\newcommand{\A}{\mathbb{A}}
\newcommand{\PP}{\mathbb{P}}
\newcommand{\E}{\mathbb{E}}
\newcommand{\eps}{\varepsilon}
\newcommand{\dd}{\,\mathrm{d}}
\begin{document}
\section*{UnsolvedMath ID: 5000001}
\textbf{Source problem code:} catalogue entry
\subsection{Definitions and mathematical statement}
Let $n\geq3$ and let $P_n$ be a regular Euclidean $n$-gon. A nonsingular billiard trajectory is a polygonal path whose interior segments lie in $P_n$ and whose reflections at edge interiors obey equal incidence and reflection angles; it does not pass through a vertex. Its length is the sum of its segment lengths. A closed trajectory returns to its initial point with its initial direction.

The source calls a trajectory preclosed if its two endpoints divide their respective boundary edges into equal-length parts and make equal angles with those oriented edges. With consistently oriented edges, a natural precise reading is equality of the oriented arclength positions along the two equal-length edges and equality of the corresponding oriented angles. The record does not specify whether endpoint subdivisions are ordered or whether incoming and outgoing angles receive different conventions. It calls closed and preclosed trajectories strongly parallel when their sequences of reflection edges agree.

For a family of parallel closed trajectories, apply the source's cancellation of all common factors in their length ratios. Compare the resulting closed-trajectory length factors with the lengths of the corresponding preclosed trajectories. The conjecture is that none of the resulting integer factors is equal to $n$. In symbolic form, if $\mathcal I_n(\mathcal F)$ denotes exactly the collection of integer factors produced by that construction for an admissible family $\mathcal F$, the target is
\[
n\notin\mathcal I_n(\mathcal F)\qquad\text{for every }n\geq3\text{ and every admissible family }\mathcal F.
\]
This notation does not supply a missing construction: the source record does not define the family, the pairing with preclosed trajectories, the common factors being cancelled, or how the final integers are extracted. In particular, the stated exclusion is equality to $n$, not an unprovided stronger divisibility assertion.
\subsection{Short English statement}
For parallel billiard trajectories in a regular $n$-gon, compare closed and preclosed lengths after the source's common-factor cancellation. Prove that the resulting integer factors never equal $n$; the exact reduction and correspondence rules are missing from the record.
\subsection{Sources}
\begin{itemize}
\item Dmitry Fuchs, \emph{Billiard Trajectories in Regular Polygons and Foliations by Closed Geodesics on Surfaces} (2020), Conjecture 1.7, \url{https://amj.math.stonybrook.edu/pdf-Springer-final/020-0170.pdf}; accessed 2026-07-29, with PDF-text extraction recorded in the supplied provenance.
\item The supplied summary mentions earlier work of Fuchs with Tabachnikov on periodic trajectories in regular polyhedra/polygons and the closed/preclosed framework, without a separate publication citation.
\item Local registry \texttt{data/unsolvedmath/problems.json}, permanent UnsolvedMath ID $5000001$, source code catalogue entry, \emph{Closed-versus-preclosed trajectory lengths}. Dataset attribution: Ulam AI and the UnsolvedMath Contributors.
\end{itemize}
\end{document}
